\documentclass[11pt,a4paper]{article}

\usepackage[utf8]{inputenc}
\usepackage[T1]{fontenc}
\usepackage[french]{babel}
\usepackage[a4paper,margin=2.2cm]{geometry}
\usepackage{amsmath,amssymb,amsthm,mathtools}
\usepackage{lmodern}
\usepackage{enumitem}
\usepackage{hyperref}
\usepackage{xcolor}
\usepackage{fancyhdr}
\usepackage{stmaryrd}
\usepackage{mathrsfs}
\usepackage{array}

\hypersetup{
    colorlinks=true,
    linkcolor=black,
    urlcolor=blue,
    citecolor=black
}

\setlength{\parindent}{0pt}
\setlength{\parskip}{0.6em}

\pagestyle{fancy}
\fancyhf{}
\fancyhead[L]{Algèbre --- Chapitre 14}
\fancyhead[R]{Polynômes}
\fancyfoot[C]{\thepage}

% Environnements
\newtheorem{definition}{Définition}[section]
\newtheorem{theoreme}[definition]{Théorème}
\newtheorem{proposition}[definition]{Proposition}
\newtheorem{corollaire}[definition]{Corollaire}
\newtheorem{lemme}[definition]{Lemme}
\theoremstyle{remark}
\newtheorem{remarque}[definition]{Remarque}
\newtheorem{exemple}[definition]{Exemple}
\newtheorem*{methode}{Méthode}
\newtheorem*{idee}{Idée}

% Commandes
\newcommand{\R}{\mathbb{R}}
\newcommand{\C}{\mathbb{C}}
\newcommand{\Q}{\mathbb{Q}}
\newcommand{\Z}{\mathbb{Z}}
\newcommand{\N}{\mathbb{N}}
\newcommand{\K}{\mathbb{K}}
\newcommand{\degP}{\mathrm{deg}}
\newcommand{\Vect}{\mathrm{Vect}}
\newcommand{\Ann}{\mathrm{Ann}}
\newcommand{\Lin}{\mathcal{L}}
\newcommand{\id}{\mathrm{Id}}
\newcommand{\Ker}{\mathrm{Ker}}
\newcommand{\Img}{\mathrm{Im}}

\begin{document}

\begin{center}
    {\LARGE \textbf{Chapitre 14 --- Polynômes}}\\[0.4em]
    {\large Structure de \(\K[X]\), racines, factorisation et préparation à la réduction}
\end{center}

\hrule
\vspace{1em}

\tableofcontents

\newpage

\section{Pourquoi ce chapitre ?}

Les polynômes interviennent déjà naturellement en algèbre élémentaire, mais leur rôle devient beaucoup plus profond en algèbre linéaire :
\begin{itemize}
    \item ils permettent de décrire les valeurs propres ;
    \item ils interviennent dans les polynômes caractéristiques et annulateurs ;
    \item ils servent à étudier les endomorphismes par \(P(u)\) ou les matrices par \(P(A)\) ;
    \item ils sont au cœur de la diagonalisation et de la réduction.
\end{itemize}

\begin{idee}
Un polynôme est à la fois un objet algébrique formel et une fonction. En algèbre linéaire, c’est surtout l’objet formel qui est fondamental.
\end{idee}

\section{L’anneau \(\K[X]\)}

\subsection{Définition}

\begin{definition}
On note \(\K[X]\) l’ensemble des polynômes à coefficients dans \(\K\), c’est-à-dire des expressions de la forme
\[
P(X)=a_0+a_1X+\cdots+a_nX^n
\]
où \(n\in\N\) et \(a_0,\dots,a_n\in\K\).
\end{definition}

\begin{remarque}
Un même polynôme peut être écrit avec des coefficients nuls supplémentaires :
\[
P(X)=a_0+a_1X+\cdots+a_nX^n+0\cdot X^{n+1}+\cdots
\]
sans changer l’objet.
\end{remarque}

\subsection{Égalité de deux polynômes}

\begin{definition}
Deux polynômes
\[
P(X)=\sum_{k=0}^n a_kX^k,
\qquad
Q(X)=\sum_{k=0}^m b_kX^k
\]
sont égaux si et seulement si tous leurs coefficients correspondants sont égaux :
\[
\forall k\in\N,\quad a_k=b_k.
\]
\end{definition}

\begin{remarque}
L’égalité de polynômes est une égalité de coefficients. Ce n’est pas seulement une égalité de valeurs prises comme fonctions, même si sur un corps infini les deux notions coïncident.
\end{remarque}

\subsection{Structure algébrique}

\begin{proposition}
L’ensemble \(\K[X]\), muni de l’addition usuelle et du produit usuel, est un anneau commutatif unitaire.
\end{proposition}

\begin{proof}
Les propriétés se vérifient coefficient par coefficient pour l’addition, et par convolution des coefficients pour le produit. L’élément neutre pour l’addition est le polynôme nul, et pour le produit le polynôme constant \(1\).
\end{proof}

\begin{remarque}
\(\K[X]\) est même un \(\K\)-espace vectoriel de dimension infinie, de base \((1,X,X^2,\dots)\).
\end{remarque}

\section{Degré d’un polynôme}

\subsection{Définition}

\begin{definition}
Soit
\[
P(X)=a_0+a_1X+\cdots+a_nX^n
\]
un polynôme non nul. Son \textbf{degré} est le plus grand entier \(k\) tel que \(a_k\neq 0\). On le note
\[
\degP(P).
\]
\end{definition}

\begin{definition}
Par convention, le degré du polynôme nul est
\[
\degP(0)=-\infty.
\]
\end{definition}

\begin{remarque}
Cette convention est très pratique, car elle permet d’écrire des formules uniformes comme
\[
\degP(PQ)=\degP(P)+\degP(Q)
\]
même lorsque l’un des polynômes est nul.
\end{remarque}

\subsection{Coefficient dominant}

\begin{definition}
Si \(P\neq 0\), le coefficient du terme de plus haut degré est appelé \textbf{coefficient dominant} de \(P\).
\end{definition}

\begin{definition}
Un polynôme non nul est dit \textbf{unitaire} si son coefficient dominant vaut \(1\).
\end{definition}

\begin{exemple}
\[
P(X)=3X^4-2X+1
\]
est de degré \(4\), de coefficient dominant \(3\), et n’est pas unitaire.
\end{exemple}

\subsection{Propriétés du degré}

\begin{proposition}
Pour tous polynômes \(P,Q\in\K[X]\),
\[
\degP(P+Q)\le \max(\degP(P),\degP(Q)).
\]
\end{proposition}

\begin{proof}
Lorsqu’on additionne les coefficients, les termes de plus haut degré peuvent éventuellement se simplifier, mais on ne peut pas créer un terme de degré supérieur au maximum initial.
\end{proof}

\begin{proposition}
Pour tous polynômes \(P,Q\in\K[X]\),
\[
\degP(PQ)=\degP(P)+\degP(Q).
\]
\end{proposition}

\begin{proof}
Si \(P\neq 0\) et \(Q\neq 0\), écrivons
\[
P=a_nX^n+\cdots,\qquad Q=b_mX^m+\cdots
\]
avec \(a_n\neq 0\) et \(b_m\neq 0\). Alors le coefficient de \(X^{n+m}\) dans \(PQ\) est
\[
a_nb_m\neq 0,
\]
puisque \(\K\) est un corps. Donc
\[
\degP(PQ)=n+m.
\]
Le cas du polynôme nul est pris en charge par la convention \(-\infty\).
\end{proof}

\begin{corollaire}
Un produit de deux polynômes non nuls n’est jamais nul.
\end{corollaire}

\begin{proof}
Si \(P\neq 0\) et \(Q\neq 0\), alors
\[
\degP(PQ)=\degP(P)+\degP(Q)\in\N,
\]
donc \(PQ\neq 0\).
\end{proof}

\section{Polynômes particuliers}

\subsection{Polynômes constants}

\begin{definition}
Les polynômes de degré \(\le 0\) sont les polynômes constants.
\end{definition}

\begin{proposition}
Les polynômes inversibles de \(\K[X]\) sont exactement les polynômes constants non nuls.
\end{proposition}

\begin{proof}
Si \(P\) est inversible, il existe \(Q\) tel que
\[
PQ=1.
\]
Donc
\[
\degP(P)+\degP(Q)=0,
\]
ce qui impose \(\degP(P)=0\). Réciproquement, tout polynôme constant non nul est inversible dans \(\K[X]\).
\end{proof}

\subsection{Polynômes de degré \(1\)}

Les polynômes de degré \(1\) sont de la forme
\[
aX+b
\qquad \text{avec } a\neq 0.
\]

Ils sont fondamentaux pour la factorisation.

\section{Évaluation}

\subsection{Valeur d’un polynôme en un scalaire}

\begin{definition}
Soit \(P(X)=a_0+a_1X+\cdots+a_nX^n\in\K[X]\), et soit \(\alpha\in\K\).  
On définit
\[
P(\alpha)=a_0+a_1\alpha+\cdots+a_n\alpha^n.
\]
\end{definition}

\begin{remarque}
L’application
\[
\mathrm{ev}_\alpha:\K[X]\to\K,\qquad P\mapsto P(\alpha)
\]
est un morphisme d’anneaux.
\end{remarque}

\begin{proposition}
Pour tous \(P,Q\in\K[X]\) et tout \(\alpha\in\K\),
\[
(P+Q)(\alpha)=P(\alpha)+Q(\alpha),
\]
\[
(PQ)(\alpha)=P(\alpha)Q(\alpha).
\]
\end{proposition}

\section{Division euclidienne}

\subsection{Énoncé}

\begin{theoreme}[Division euclidienne]
Soient \(A,B\in\K[X]\), avec \(B\neq 0\).  
Il existe un unique couple \((Q,R)\in\K[X]^2\) tel que
\[
A=BQ+R
\]
avec
\[
\degP(R)<\degP(B).
\]
\end{theoreme}

\subsection{Démonstration : unicité}

\begin{proof}[Démonstration de l’unicité]
Supposons
\[
A=BQ+R=BQ'+R'
\]
avec
\[
\degP(R)<\degP(B),
\qquad
\degP(R')<\degP(B).
\]
Alors
\[
B(Q-Q')=R'-R.
\]
Si \(Q-Q'\neq 0\), alors
\[
\degP(B(Q-Q'))\ge \degP(B),
\]
tandis que
\[
\degP(R'-R)<\degP(B),
\]
contradiction. Donc \(Q=Q'\), puis \(R=R'\).
\end{proof}

\subsection{Démonstration : existence}

\begin{proof}[Idée de la démonstration de l’existence]
On raisonne par récurrence sur le degré de \(A\).  
Si
\[
\degP(A)<\degP(B),
\]
on prend \(Q=0\) et \(R=A\).

Sinon, si
\[
A=a_nX^n+\cdots,\qquad B=b_mX^m+\cdots
\]
avec \(n\ge m\), on soustrait à \(A\) le polynôme
\[
\frac{a_n}{b_m}X^{n-m}B,
\]
ce qui fait disparaître le terme dominant. On recommence ensuite.
\end{proof}

\subsection{Cas particulier : division par \(X-\alpha\)}

\begin{proposition}
Soit \(P\in\K[X]\) et \(\alpha\in\K\). Alors il existe un unique couple \((Q,r)\in\K[X]\times \K\) tel que
\[
P(X)=(X-\alpha)Q(X)+r.
\]
De plus,
\[
r=P(\alpha).
\]
\end{proposition}

\begin{proof}
Le reste dans la division par un polynôme de degré \(1\) est de degré strictement inférieur à \(1\), donc constant. En évaluant en \(X=\alpha\), on obtient
\[
P(\alpha)=r.
\]
\end{proof}

\begin{corollaire}[Théorème des racines]
\[
P(\alpha)=0
\Longleftrightarrow
X-\alpha \text{ divise } P.
\]
\end{corollaire}

\section{Racines d’un polynôme}

\subsection{Définition}

\begin{definition}
Un scalaire \(\alpha\in\K\) est une \textbf{racine} de \(P\in\K[X]\) si
\[
P(\alpha)=0.
\]
\end{definition}

\begin{exemple}
Les racines de
\[
P(X)=X^2-1
\]
dans \(\R\) sont \(1\) et \(-1\).
\end{exemple}

\subsection{Nombre de racines}

\begin{theoreme}
Un polynôme non nul de degré \(n\) admet au plus \(n\) racines distinctes dans \(\K\).
\end{theoreme}

\begin{proof}
Par récurrence sur \(n\).

Si \(n=0\), le polynôme est constant non nul et n’a pas de racine.

Supposons le résultat vrai jusqu’au rang \(n-1\). Soit \(P\) de degré \(n\). Si \(P\) n’a pas de racine, c’est terminé. Sinon, soit \(\alpha\) une racine. Alors
\[
P=(X-\alpha)Q
\]
avec
\[
\degP(Q)=n-1.
\]
Toute racine distincte de \(\alpha\) est racine de \(Q\). Par hypothèse de récurrence, \(Q\) a au plus \(n-1\) racines distinctes. Donc \(P\) a au plus \(n\) racines distinctes.
\end{proof}

\begin{corollaire}
Si deux polynômes de degré \(\le n\) coïncident en \(n+1\) points distincts, alors ils sont égaux.
\end{corollaire}

\begin{proof}
Leur différence est un polynôme de degré \(\le n\) ayant \(n+1\) racines distinctes, donc c’est le polynôme nul.
\end{proof}

\section{Multiplicité d’une racine}

\subsection{Définition}

\begin{definition}
Soit \(\alpha\in\K\). On dit que \(\alpha\) est racine de multiplicité \(m\ge 1\) de \(P\) si
\[
P(X)=(X-\alpha)^mQ(X)
\]
avec
\[
Q(\alpha)\neq 0.
\]
\end{definition}

\begin{remarque}
La multiplicité est le plus grand entier \(m\) tel que \((X-\alpha)^m\) divise \(P\).
\end{remarque}

\begin{exemple}
Dans
\[
P(X)=(X-1)^3(X+2)^2,
\]
la racine \(1\) est de multiplicité \(3\), et la racine \(-2\) est de multiplicité \(2\).
\end{exemple}

\subsection{Somme des multiplicités}

\begin{proposition}
La somme des multiplicités des racines distinctes d’un polynôme non nul est inférieure ou égale à son degré.
\end{proposition}

\begin{proof}
Si
\[
P(X)=\prod_{i=1}^r (X-\alpha_i)^{m_i}Q(X)
\]
avec \(Q(\alpha_i)\neq 0\), alors
\[
\degP(P)\ge \sum_{i=1}^r m_i.
\]
\end{proof}

\section{Dérivée formelle d’un polynôme}

\subsection{Définition}

\begin{definition}
Si
\[
P(X)=a_0+a_1X+\cdots+a_nX^n,
\]
on définit sa \textbf{dérivée formelle} par
\[
P'(X)=a_1+2a_2X+\cdots+na_nX^{n-1}.
\]
\end{definition}

\begin{remarque}
C’est la même formule que pour la dérivation usuelle, mais ici elle est définie algébriquement dans \(\K[X]\).
\end{remarque}

\subsection{Propriétés}

\begin{proposition}
Pour tous \(P,Q\in\K[X]\) et \(\lambda\in\K\),
\[
(P+Q)'=P'+Q',
\qquad
(\lambda P)'=\lambda P',
\qquad
(PQ)'=P'Q+PQ'.
\]
\end{proposition}

\begin{proof}
Les deux premières propriétés se vérifient coefficient par coefficient.  
Pour le produit, on développe puis on regroupe les coefficients.
\end{proof}

\subsection{Lien avec les racines multiples}

\begin{theoreme}
Soit \(\alpha\in\K\). Alors \(\alpha\) est racine multiple de \(P\) si et seulement si
\[
P(\alpha)=0
\qquad \text{et} \qquad
P'(\alpha)=0.
\]
\end{theoreme}

\begin{proof}
Supposons
\[
P(X)=(X-\alpha)^mQ(X)
\]
avec \(m\ge 2\). Alors
\[
P'(X)=m(X-\alpha)^{m-1}Q(X)+(X-\alpha)^mQ'(X),
\]
donc \(P'(\alpha)=0\).

Réciproquement, si \(P(\alpha)=0\), alors
\[
P(X)=(X-\alpha)Q(X).
\]
En dérivant :
\[
P'(X)=Q(X)+(X-\alpha)Q'(X),
\]
donc
\[
P'(\alpha)=Q(\alpha).
\]
Si \(P'(\alpha)=0\), alors \(Q(\alpha)=0\), donc \(X-\alpha\) divise \(Q\), et par suite \((X-\alpha)^2\) divise \(P\).
\end{proof}

\begin{corollaire}
\(\alpha\) est racine de multiplicité au moins \(m\) de \(P\) si et seulement si
\[
P(\alpha)=P'(\alpha)=\cdots=P^{(m-1)}(\alpha)=0.
\]
\end{corollaire}

\section{Factorisation}

\subsection{Sur \(\C\)}

\begin{theoreme}[Théorème fondamental de l’algèbre]
Tout polynôme non constant de \(\C[X]\) admet au moins une racine complexe.
\end{theoreme}

\begin{corollaire}
Tout polynôme non constant de \(\C[X]\) est scindé, c’est-à-dire se factorise sous la forme
\[
P(X)=a\prod_{k=1}^r (X-\alpha_k)^{m_k}.
\]
\end{corollaire}

\begin{proof}
On itère le théorème des racines : à chaque racine \(\alpha\), on factorise par \(X-\alpha\), puis on recommence sur le quotient.
\end{proof}

\subsection{Sur \(\R\)}

\begin{proposition}
Tout polynôme non constant de \(\R[X]\) se factorise en produit de polynômes irréductibles de degré \(1\) ou \(2\).
\end{proposition}

\begin{proof}
Dans \(\C[X]\), le polynôme est scindé. Les racines complexes non réelles apparaissent par paires conjuguées, et chaque paire
\[
(X-z)(X-\overline z)
\]
donne un polynôme réel irréductible de degré \(2\).
\end{proof}

\begin{exemple}
\[
X^4+1
\]
n’est pas scindé sur \(\R\), mais il se factorise en produit de deux polynômes irréductibles de degré \(2\).
\end{exemple}

\section{Polynômes irréductibles}

\subsection{Définition}

\begin{definition}
Un polynôme \(P\in\K[X]\) est dit \textbf{irréductible} s’il est non constant et s’il ne peut pas s’écrire comme produit de deux polynômes de degrés strictement positifs.
\end{definition}

\begin{exemple}
Dans \(\R[X]\), le polynôme
\[
X^2+1
\]
est irréductible.
\end{exemple}

\begin{exemple}
Dans \(\C[X]\), les polynômes irréductibles sont exactement les polynômes de degré \(1\).
\end{exemple}

\subsection{Cas du degré \(2\) ou \(3\)}

\begin{proposition}
Un polynôme réel de degré \(2\) ou \(3\) est irréductible sur \(\R\) si et seulement s’il n’a pas de racine réelle.
\end{proposition}

\begin{proof}
S’il a une racine réelle \(\alpha\), il est divisible par \(X-\alpha\), donc réductible.

Réciproquement, s’il est réductible, il s’écrit comme produit de deux polynômes non constants. Comme le degré est \(2\) ou \(3\), l’un des deux facteurs est nécessairement de degré \(1\), donc fournit une racine réelle.
\end{proof}

\section{PGCD et polynômes premiers entre eux}

\subsection{Définition}

\begin{definition}
Un polynôme \(D\in\K[X]\) est un \textbf{plus grand commun diviseur} de \(P\) et \(Q\) s’il vérifie :
\begin{enumerate}[label=\arabic*)]
    \item \(D\) divise \(P\) et \(Q\) ;
    \item tout diviseur commun de \(P\) et \(Q\) divise \(D\).
\end{enumerate}
\end{definition}

\begin{remarque}
Le PGCD est défini à multiplication près par une constante non nulle. On le choisit souvent unitaire.
\end{remarque}

\subsection{Bézout}

\begin{theoreme}[Bézout]
Deux polynômes \(P,Q\in\K[X]\) sont premiers entre eux si et seulement s’il existe \(U,V\in\K[X]\) tels que
\[
UP+VQ=1.
\]
\end{theoreme}

\begin{remarque}
C’est un outil essentiel pour la suite, notamment dans les décompositions de noyaux et les polynômes annulateurs.
\end{remarque}

\section{Polynômes annulateurs}

\subsection{Définition}

\begin{definition}
Soit \(u\in\Lin(E,E)\) un endomorphisme.  
Un polynôme \(P\in\K[X]\) est appelé \textbf{polynôme annulateur} de \(u\) si
\[
P(u)=0.
\]
\end{definition}

De même, pour une matrice \(A\in M_n(\K)\), un polynôme \(P\) est annulateur si
\[
P(A)=0.
\]

\subsection{Lien avec les valeurs propres}

\begin{proposition}
Si \(P\) est un polynôme annulateur de \(u\), alors toute valeur propre \(\lambda\) de \(u\) vérifie
\[
P(\lambda)=0.
\]
\end{proposition}

\begin{proof}
Si \(x\neq 0\) est vecteur propre associé à \(\lambda\), alors
\[
u(x)=\lambda x.
\]
On en déduit
\[
P(u)(x)=P(\lambda)x.
\]
Comme \(P(u)=0\), on a
\[
P(\lambda)x=0.
\]
Et comme \(x\neq 0\),
\[
P(\lambda)=0.
\]
\end{proof}

\begin{remarque}
Ce résultat sera central dans les chapitres de réduction.
\end{remarque}

\subsection{Polynôme minimal}

\begin{definition}
Le \textbf{polynôme minimal} d’un endomorphisme \(u\) est le polynôme unitaire de plus petit degré qui annule \(u\).
\end{definition}

\begin{remarque}
Même si son étude détaillée peut être repoussée, il est utile de le mentionner dès maintenant car il structure une grande partie de la réduction.
\end{remarque}

\section{Polynômes scindés}

\begin{definition}
Un polynôme \(P\in\K[X]\) est dit \textbf{scindé sur \(\K\)} s’il se factorise en produit de facteurs de degré \(1\) :
\[
P(X)=a\prod_{i=1}^r (X-\alpha_i)^{m_i}.
\]
\end{definition}

\begin{remarque}
Sur \(\C\), tout polynôme non constant est scindé.  
Sur \(\R\), ce n’est pas toujours le cas.
\end{remarque}

\section{Fractions rationnelles}

\subsection{Définition}

\begin{definition}
Une \textbf{fraction rationnelle} sur \(\K\) est une expression de la forme
\[
\frac{P(X)}{Q(X)}
\]
où \(P,Q\in\K[X]\) et \(Q\neq 0\).
\end{definition}

\begin{remarque}
Deux fractions rationnelles
\[
\frac{P}{Q}
\qquad \text{et} \qquad
\frac{P'}{Q'}
\]
sont égales si
\[
PQ'=P'Q.
\]
\end{remarque}

\subsection{Décomposition en éléments simples (aperçu)}

\begin{remarque}
Dans de nombreuses filières, les fractions rationnelles sont surtout utiles pour l’intégration ou les suites récurrentes linéaires. On peut les décomposer en éléments simples lorsque le dénominateur est factorisé.
\end{remarque}

\begin{exemple}
Sur \(\R\),
\[
\frac{1}{X(X-1)}=-\frac{1}{X}+\frac{1}{X-1}.
\]
\end{exemple}

\section{Méthodes pratiques}

\begin{methode}
Pour montrer qu’un scalaire \(\alpha\) est racine d’un polynôme \(P\), il suffit de calculer \(P(\alpha)\) et de vérifier que
\[
P(\alpha)=0.
\]
\end{methode}

\begin{methode}
Pour factoriser un polynôme :
\begin{enumerate}
    \item chercher des racines évidentes ;
    \item utiliser le théorème des racines ;
    \item effectuer une division euclidienne ;
    \item recommencer sur le quotient ;
    \item selon le corps de base, conclure par une factorisation complète ou partielle.
\end{enumerate}
\end{methode}

\begin{methode}
Pour détecter une racine multiple, vérifier si
\[
P(\alpha)=0
\qquad \text{et} \qquad
P'(\alpha)=0.
\]
\end{methode}

\begin{methode}
Pour étudier un endomorphisme à l’aide d’un polynôme annulateur, commencer par repérer les racines du polynôme : elles contiennent toutes les valeurs propres possibles.
\end{methode}

\section{Erreurs classiques}

\begin{itemize}
    \item Confondre le polynôme nul et le polynôme constant \(0\) sans utiliser correctement la convention sur le degré.
    \item Croire qu’un polynôme de degré \(n\) a toujours exactement \(n\) racines dans \(\K\).
    \item Oublier que les racines complexes non réelles d’un polynôme réel apparaissent par paires conjuguées.
    \item Confondre racine simple et racine multiple.
    \item Passer trop vite d’une égalité de fonctions à une égalité de polynômes sans justifier le cadre.
    \item Oublier que la division euclidienne n’existe que par un polynôme non nul.
\end{itemize}

\section{Résumé du chapitre}

\begin{itemize}
    \item \(\K[X]\) est un anneau commutatif unitaire et un \(\K\)-espace vectoriel.
    \item Le degré mesure le terme dominant d’un polynôme non nul.
    \item On a
    \[
    \degP(PQ)=\degP(P)+\degP(Q).
    \]
    \item Le théorème de division euclidienne permet d’écrire
    \[
    A=BQ+R
    \quad \text{avec} \quad
    \degP(R)<\degP(B).
    \]
    \item \(\alpha\) est racine de \(P\) si et seulement si \(X-\alpha\) divise \(P\).
    \item Un polynôme non nul de degré \(n\) a au plus \(n\) racines distinctes.
    \item Une racine est multiple si et seulement si elle annule aussi la dérivée formelle.
    \item Sur \(\C\), tout polynôme est scindé.
    \item Les polynômes annulateurs joueront un rôle central dans la réduction.
\end{itemize}

\section*{Exercices d’application immédiate}

\begin{enumerate}[label=\textbf{Exercice \arabic*.}]
    \item Déterminer le degré et le coefficient dominant de
    \[
    P(X)=3X^4-2X+1.
    \]

    \item Effectuer la division euclidienne de
    \[
    X^3-2X+1
    \]
    par
    \[
    X-1.
    \]

    \item Montrer que \(1\) est racine de
    \[
    X^3-2X+1.
    \]

    \item Factoriser
    \[
    X^3-2X+1
    \]
    dans \(\R[X]\).

    \item Déterminer les racines et leurs multiplicités de
    \[
    (X-1)^3(X+2)^2.
    \]

    \item Montrer que \(X^2+1\) est irréductible dans \(\R[X]\).

    \item Montrer que \(1\) est racine multiple de
    \[
    X^3-3X^2+3X-1.
    \]

    \item Calculer la dérivée formelle de
    \[
    P(X)=X^5-4X^3+2X-7.
    \]

    \item Montrer qu’un polynôme de degré \(n\) ne peut pas avoir plus de \(n\) racines distinctes.

    \item Soit \(u\) un endomorphisme tel que
    \[
    u^2-3u+2\id=0.
    \]
    Montrer que ses valeurs propres éventuelles sont parmi \(1\) et \(2\).
\end{enumerate}

\section{Interpolation de Lagrange}

\begin{theoreme}
Soient \(a_0,\dots,a_n\) des scalaires deux à deux distincts.
L’application
\[
u:\K_n[X]\to\K^{n+1},\qquad P\mapsto (P(a_0),\dots,P(a_n))
\]
est un isomorphisme.
\end{theoreme}

\begin{proof}
Si \(P\in\K_n[X]\) s’annule en \(n+1\) points distincts, alors \(P\) possède au moins \(n+1\) racines ; comme \(\deg(P)\le n\), on obtient \(P=0\).
L’application est donc injective. Les espaces de départ et d’arrivée ayant même dimension \(n+1\), elle est bijective.
\end{proof}

\begin{definition}
Pour \(i\in\{0,\dots,n\}\), le \textbf{polynôme de Lagrange} associé à \(a_i\) est
\[
L_i(X)=\prod_{j\ne i}\frac{X-a_j}{a_i-a_j}.
\]
Il vérifie \(L_i(a_j)=\delta_{ij}\).
\end{definition}

\begin{proposition}
La famille \((L_0,\dots,L_n)\) est une base de \(\K_n[X]\), et pour tout \(P\in\K_n[X]\),
\[
P(X)=\sum_{i=0}^n P(a_i)L_i(X).
\]
\end{proposition}

\end{document}
